% =====================================================================================
% Contribution au document du Lab Cyber envoye par Hugo Rapior le 2 septembre 2026 a 10h48.
% Deux sections : « Rayonnement et credibilite », completee, et « Vision 2030 », ecrite.
% CONTRAINTE : deux pages au maximum.
%
% VERSION 2, DU MEME JOUR. La premiere version reprenait deux dates de la note interne qui ne
% resistent pas a la verification, et sourcait le reste sur des resumes de recherche. Cinq pages
% sources ont ete lues depuis. Ce qui a change :
%   - « 2 aout 2026, obligations haut risque de l'IA Act » est FAUX : l'annexe III est reportee
%     au 2 decembre 2027. Ce qui reste au 2 aout 2026 est la transparence (art. 50) et la
%     litteratie (art. 4) ;
%   - « 17 janvier 2027, DORA aux petites entites » ne correspond a AUCUNE echeance documentee.
%     Remplace par les deux jalons reels du CRA ;
%   - le nombre de prestataires critiques et leurs noms ne sont PAS dans le communique des ESA :
%     ils sont donnes avec leur statut de preuve, pas comme un fait etabli ;
%   - l'article 58 est cite au verbatim, sans la mention de proposition legislative qui venait
%     d'un resume et non du texte ;
%   - la strategie nationale est attribuee au SGDSN et non a l'ANSSI, et la citation sur la
%     « guerre de haute intensite » est retiree faute de source.
%
% CE QUE « COMPLETER » VEUT DIRE ICI. La liste est celle d'Hugo, dans SON document : il sait ce
% qui s'est passe. On ne la corrige pas, on l'AUGMENTE de ce qu'elle omet.
%
% Aucun tiret cadratin : le babel francais en met un en puce d'itemize par defaut.
% =====================================================================================
\documentclass[11pt,a4paper]{article}

\usepackage[utf8]{inputenc}
\usepackage[T1]{fontenc}
\usepackage{lmodern}
\usepackage[french]{babel}
\usepackage[top=1.4cm,bottom=1.3cm,left=1.8cm,right=1.8cm]{geometry}
\usepackage{array}
\usepackage{enumitem}
\usepackage{xcolor}
\usepackage{eurosym}
\usepackage{titlesec}

\definecolor{encre}{HTML}{1B1E30}
\definecolor{accent}{HTML}{A6002E}
\definecolor{bleu}{HTML}{204993}
\definecolor{grisclair}{HTML}{F2F2F2}

\color{encre}
\titleformat{\section}{\normalfont\large\bfseries\color{bleu}}{}{0pt}{}
\titlespacing{\section}{0pt}{9pt}{3pt}
\titleformat{\subsection}{\normalfont\bfseries}{}{0pt}{}
\titlespacing{\subsection}{0pt}{7pt}{1pt}
\setlist[itemize]{label=\textcolor{accent}{\textbullet},itemsep=2pt,topsep=2pt,
                  leftmargin=1.1em,parsep=0pt}
\newcommand{\src}[1]{{\footnotesize\itshape\textcolor{bleu}{Source : #1}}}
\newcommand{\encadre}[2]{%
  \vspace{2pt}\noindent\colorbox{grisclair}{%
  \begin{minipage}{\dimexpr\textwidth-2\fboxsep}
  \footnotesize\textbf{#1} #2
  \end{minipage}}\vspace{3pt}}

\pagestyle{plain}

\begin{document}
\small

{\Large\bfseries\color{bleu} Lab Cyber : rayonnement et vision 2030}\\[3pt]
{\color{accent}\rule{\textwidth}{1.1pt}}\\[3pt]
{\footnotesize Contribution de Kélian Kaddouri \hfill Pour Hugo Rapior \textbullet{} 2 septembre 2026}

\vspace{4pt}

\section{Rayonnement et crédibilité}

\noindent Ta liste est reprise telle quelle, avec la précision factuelle là où elle rend la
ligne plus solide, puis trois livrables qui n'y figuraient pas.

\begin{itemize}
\item une vidéo de présentation publiée sur LinkedIn ;
\item la participation au livre blanc FCC $\times$ Finance Innovation, \emph{Fondations et
      gouvernance pour une IA de confiance} : qualité de la donnée, dispositifs de gouvernance
      et de contrôle, cas d'usage sur les trois lignes de défense, parangonnage des solutions
      d'IA déployées en banque et en assurance dont la fraude et la cybersécurité ;
\item les interventions dans des formations actuarielles ;
\item la présentation au European Congress of Actuaries, qui se tient à Paris les 18 et 19 juin
      2026 et dont l'appel à intervenants comportait \og{}Role of actuaries in modelling Cyber
      Risk\fg{} au titre des risques émergents ;
\item les articles et publications sur l'IA Act, les risques géopolitiques, le risque cyber et
      l'IA de confiance, dont l'article du 10 juillet 2025 sur le \emph{GPAI Code of Practice},
      adossé à l'article 53 du règlement sur l'IA et à ses annexes XI et XII ;
\item les travaux ANSSI et PANAME sur l'audit des modèles d'IA. PANAME a été lancé en juin 2025
      par l'ANSSI, la CNIL, le PEReN et le projet IPoP du PEPR Cybersécurité piloté par Inria,
      et son appel à manifestation d'intérêt s'est tenu du 26 février au 28 mars 2026.
\end{itemize}

\subsection{Trois livrables à ajouter, et le premier est le plus visible de tous}

\begin{itemize}
\item \textbf{Les rapports LUCY 2025 et LUCY 2026}, analyse du marché de la cyberassurance en
      France. Assis sur 14\,124 polices et 448 sinistres déclarés auprès d'Aon, Marsh et WTW :
      ratio de sinistralité de 17 \% en 2024 contre 12 \% en 2023, nombre de sinistres en baisse
      de 27 \% pour des montants indemnisés en hausse de 43 \%, retour des sinistres XXL avec
      21 M\euro{} sur les incidents majeurs, tarifs en repli de 19,8 \% sur les grandes
      entreprises et de 8,7 \% sur les ETI, capacités en hausse de 23 \%, franchises en baisse
      de 15 \% pour les grands comptes contre une hausse de 231 \% pour les PME.
      \textbf{C'est le livrable cyber du Lab le plus cité à l'extérieur}, et il ne figurait pas
      dans la liste.
\item \textbf{L'atelier du 24 novembre 2025} à la journée 100 \% Actuaires, 100 \% Data Science
      \& Durabilité : \og{}Scénarios stochastiques, données textuelles et IA pour une meilleure
      quantification du risque cyber\fg{}. Sa valeur de rayonnement tient autant à son contenu
      qu'à sa \emph{composition} : une co-intervention avec un assureur cyber et une direction de
      master académique. C'est une caution externe, pas une prise de parole isolée.
\item \textbf{Un mémoire d'actuariat en cours sur la quantification du besoin de capital au
      titre de DORA}, soutenance 2026. Il alimente l'axe SCR cyber et RWA cyber, aujourd'hui le
      seul des axes annoncés du Lab qui n'ait pas encore de publication.
\end{itemize}

\noindent Le projet de podcast consacré au rapport LUCY peut constituer un support
supplémentaire de valorisation et de prospection. Sa force propre est qu'il porte sur le
livrable le plus reconnu du Lab plutôt que sur un sujet neuf : il capitalise au lieu d'ouvrir un
front.

\encadre{Une suggestion de présentation, et non une réserve sur les faits.}{Devant un client, il
vaut mieux distinguer dans cette liste ce qui laisse une \textbf{trace publique vérifiable} (le
livre blanc, les articles, les rapports LUCY) de ce qui relève d'une \textbf{intervention ou
d'une contribution} (formations, congrès, travaux de veille). Ce n'est pas une hiérarchie de
valeur mais d'usage : la première catégorie se vérifie en un clic, la seconde gagne à être
attribuée précisément. J'ai fait le relevé sur sources publiques, et c'est ce que trouve un
interlocuteur qui vérifie.}

\section{Vision 2030 : du document à la preuve}

\noindent \textbf{Un seul mouvement commande la décennie.} La demande passe de la documentation
à la preuve : registres vérifiés, tests pilotés par la menace, audits de modèles effectifs
plutôt que procédures écrites. L'écart entre l'exigence et la maturité réelle est mesuré. À
l'exercice à blanc de 2024 sur le registre d'information, \textbf{6,5 \%} des registres
analysés ont passé la totalité des \textbf{116} contrôles qualité, sur près de \textbf{mille}
entités financières de l'Union. \\
\src{ESA (EBA, EIOPA, ESMA), communiqué du 17 décembre 2024. Précision d'usage : les ESA
qualifient cette qualité de \og{}in line with the ESAs' expectations\fg{} vu le caractère
volontaire de l'exercice, et jugeaient l'objectif de 2025 \og{}not out of reach, subject to some
additional efforts from the industry\fg{}. Le 6,5 \% mesure un point de départ, pas un naufrage.}

\subsection{2 août 2026 : IA Act, transparence et littératie}
Les obligations de transparence de l'article 50 et le devoir de littératie en IA de l'article 4
s'appliquent, et restent sur leur calendrier d'origine. \\
\src{règlement (UE) 2024/1689, articles 4, 50 et 113. Le commentaire de l'accord de report
souligne que \og{}2 August 2026 remains an active compliance date\fg{}.}

\subsection{11 septembre 2026 : Cyber Resilience Act, notification des vulnérabilités}
Les fabricants de produits connectés doivent notifier les vulnérabilités activement exploitées
et les incidents de sécurité. \\
\src{Commission européenne, page officielle du CRA, entré en vigueur le 10 décembre 2024.}

\subsection{2026-2027 : supervision directe des prestataires TIC critiques}
Les autorités européennes ont désigné les premiers prestataires TIC critiques le
\textbf{18 novembre 2025}, qui passent sous supervision directe : \og{}through direct oversight
engagement, the ESAs will assess whether CTPPs have appropriate risk management and governance
frameworks in place\fg{}. \textbf{C'est le jalon le plus exploitable pour le Lab} : une liste
nominative de points de concentration transforme l'accumulation par cause commune en exercice
mesurable sur des tiers identifiés, là où elle est aujourd'hui un facteur posé. \\
\src{communiqué conjoint EBA, EIOPA et ESMA du 18 novembre 2025. \textbf{Réserve} : le
communiqué ne donne ni le nombre ni les noms, la liste figurant dans une annexe distincte. Les
commentaires publiés avancent dix-neuf prestataires et ne concordent pas sur les entreprises :
relever la liste sur l'annexe des ESA avant de nommer qui que ce soit.}

\subsection{2 décembre 2027 : IA Act, obligations haut risque de l'annexe III}
Reportées du 2 août 2026, et au 2 août 2028 pour l'IA embarquée dans les produits réglementés de
l'annexe I. Les deux points qui concernent la finance sont l'annexe III 5(b), évaluation de
solvabilité et scoring de crédit, et 5(c), évaluation du risque et tarification en assurance vie
et santé. \\
\src{accord politique provisoire des colégislateurs du 6 mai 2026, confirmé par les États
membres le 13 mai. \textbf{Réserve} : les nouvelles dates ne lient qu'après publication au
Journal officiel, attendue avant le 2 août 2026 et \textbf{que je n'ai pas vérifiée}. À
confirmer au Journal officiel avant de présenter le report comme définitif.}

\subsection{11 décembre 2027 : Cyber Resilience Act, obligations principales}
Exigences essentielles de l'annexe I, évaluation de conformité, documentation technique,
déclaration UE de conformité et marquage CE. \\
\src{Commission européenne, page officielle du CRA.}

\subsection{17 janvier 2028 : réexamen de DORA, et c'est le jalon qui nous concerne le plus}
\og{}Au plus tard le 17 janvier 2028, la Commission, après avoir consulté les AES et le CERS,
selon le cas, procède à un réexamen et remet au Parlement européen et au Conseil un rapport\fg{}.
Deux des cinq objets du réexamen sont directement notre terrain : les \textbf{critères de
désignation des prestataires tiers critiques} et le \textbf{caractère volontaire de la
notification des cybermenaces importantes} de l'article 19. Autrement dit, le texte a déjà
inscrit à son ordre du jour la question du caractère obligatoire de la notification : une
recommandation de reporting sur ce point ne vise pas un hypothétique. C'est la seule fenêtre où
une méthode publiée peut peser sur le cadre plutôt que le subir, et elle se prépare deux ans à
l'avance. \\
\src{règlement (UE) 2022/2554, article 58, clause de réexamen, cité au verbatim.}

\subsection{2026-2030 : stratégie nationale de cybersécurité}
Publiée le \textbf{29 janvier 2026}, élaborée \textbf{sous l'égide du SGDSN} à la demande du
Président de la République, avec l'ensemble des ministères et un panel d'experts industriels,
scientifiques et académiques. Elle prolonge la Revue nationale stratégique. \\
\src{SGDSN, Stratégie nationale de cybersécurité 2026-2030. \textbf{Attention à l'attribution} :
la stratégie est pilotée par le SGDSN, l'ANSSI en étant un opérateur et non l'auteur.}

\subsection{Ce que cela dessine comme offre, et le risque à porter ouvertement}

\begin{itemize}
\item \textbf{Quatre briques, dans cet ordre de maturité} : conformité outillée aujourd'hui
      (registre, cartographie, KPI et KRI, scoring des tiers), quantification du capital en
      2026-2027 (SCR cyber, RWA cyber, pilier 2), accumulation et transfert en 2027-2028 dans le
      prolongement de LUCY, puis audit de modèles cyber et IA confondus vers 2030. La première
      finance le reste mais se banalise le plus vite : dès que l'exercice se répète, il passe à
      l'outil.
\item \textbf{Un levier non utilisé.} Le pôle mène des thèses CIFRE sur la finance durable et le
      climat, et aucune sur le cyber alors que c'est l'un des trois labs. Une CIFRE cyber
      achèterait trois ans de recherche financée et un titre publiable, ce qu'aucune mission ne
      produit.
\item \textbf{Le risque principal, et il vaut mieux l'énoncer.} Il n'existe aucun module DORA en
      formule standard : le besoin de capital cyber reste un exercice d'ORSA, donc un récit que
      l'entité \emph{choisit} de produire. Si le réexamen de 2028 ne crée pas d'exigence
      chiffrée, la quantification restera une prestation de conviction plutôt que de conformité.
      C'est la raison de préparer ce jalon plutôt que de l'attendre.
\end{itemize}

\vspace{4pt}
\noindent{\color{accent}\rule{\textwidth}{0.6pt}}\\[2pt]
{\footnotesize\itshape Sources lues directement : communiqués EBA du 17 décembre 2024 et EIOPA
du 18 novembre 2025, page officielle du CRA de la Commission, article 58 du règlement (UE)
2022/2554, page du SGDSN. Éléments Nexialog relevés sur les pages publiques du cabinet. Trois
points à confirmer avant diffusion externe : la publication au Journal officiel du report de
l'annexe III, le nombre et les noms des prestataires critiques dans l'annexe des ESA, et la
mention d'une proposition législative à l'article 58, absente du verbatim obtenu.
L'ordonnancement des briques est une hypothèse de lecture, à confronter au carnet de commandes.}

\end{document}
